\documentclass{article}
%  ╔══════════════════════════════════════════════════════════╗
%  ║                         Title                            ║
%  ╚══════════════════════════════════════════════════════════╝
%NOTE: I want \title,\author to be near top, and this \renewcommand{\maketitle} fails if these are above it.
\usepackage[utf8]{inputenc}         % Used to compile any UTF-8 Character (and supports Unicode)
\usepackage{titling}

\renewcommand{\maketitle}{
    \quad
    \vspace{1em}
  \begin{center}
      \LARGE\bfseries \thetitle \par
      \vspace{0.6em}
      \large
  \end{center}
    \vspace{1em}
}

\title{Alien Perspective on Analysis}
\author{Nathanael Chwojko-Srawley}
\date{\today}

%  ╔══════════════════════════════════════════════════════════╗
%  ║                         packages                         ║
%  ╚══════════════════════════════════════════════════════════╝

\usepackage{extarrows}         % Used to compile any UTF-8 Character (and supports Unicode)
\usepackage{stackengine}
\usepackage{colortbl}
\usepackage{scalerel}
\usepackage{amsmath, amssymb} % Used to write better mathematical expressions (ex. \mathbb{R})
\usepackage{stackrel}         % for left-right arrows, staking symbol above and bellow
\usepackage{mathrsfs}
\usepackage{bbm}                     % for mathbbm{1}
\usepackage{euscript}
\usepackage{commutative-diagrams}
\usepackage[font=itshape]{quoting}
\usepackage{mathtools}
  \usepackage{enumitem}
\usepackage{amsthm}           % Used to create theorem enviroments.
\usepackage{graphicx}         % Let's you use images in LaTeX; ex the \includegraphics command.
\usepackage{hhline}
%\usepackage{luatexja}
%\usepackage{fontspec}
\usepackage{dsfont}
\usepackage{wasysym}          % emoji's!
\usepackage{float}            % package with ``H'' for figures and tables, ex. Images, tables, etc
%\usepackage{subcaption}      % More options with sub-captions.
\usepackage{setspace}         % More flexibility on spacing in document.
\usepackage{hyperref}         % let's me put hyperlinks
%\usepackage{tikz}             % for commutative diagrams https://tex.stackexchange.com/questions/419221/how-to-do-the-pushout-with-universal-property
\usepackage{tikz-cd}          % for commutative diagrams https://tex.stackexchange.com/questions/419221/how-to-do-the-pushout-with-universal-property
%\usepackage{quiver}           % for better commutative diagrams
\usepackage{bookmark}         % creating heading shortcut bar on the left as in word
\usepackage{longtable}        % create tables that extend past one page
%\usepackage{siunitx}         % required for alignment
%\usepackage{multirow}        % for multiple rows
\usepackage{microtype}        % makes nice text. Fixes some under-/over- filled box problems.
%\usepackage{fancyvrb}        % Let's you write code (different style than listings).
\usepackage[most]{tcolorbox} 	      % Makes nice boxes for thms, cor, prop, or anything else.
%\usepackage{enumitem}        % More flexibility on enumeration (don't know much).
\usepackage{xcolor}           % Colour text.
\usepackage{wrapfig}          % To wrap text around figure
\usepackage{listings}         % Create nice colour code blocks (customized in preamble)
\usepackage{thmtools}         % Customize theorem environment (in preamble).
%\usepackage{marvosym}        % Provide ``basic'' symbols (hazard, religious, smiley face etc.)
\usepackage{array}            % \begin{table}{>{\em}c c} -> 1st column italic (bfseries for bold)
%\usepackage[english]{babel}  % If I want to blindtext in english.
\usepackage{blindtext}
\usepackage[a4paper, total={6in, 8.5in}]{geometry}
%\usepackage[symbol]{footmisc} % For daggers and asterisk for footnotes. 1: *, 2: dagger, etc.
\usepackage{fancyhdr}
\usepackage{titlesec}
%\usepackage{pst-plot}          %To make nice plots: https://tex.stackexchange.com/questions/47594/how-can-i-make-a-graph-of-a-function#47596
% \usepackage[answerdelayed]{exercise}
\usepackage{etoolbox}          % for Dynkin diagrams
\usepackage{dynkin-diagrams}   % for Dynkin diagrams

% ╔═════════════════════════════════════════════════════╗
% ║           for figures via inkspace                  ║
% ╚═════════════════════════════════════════════════════╝

% to import figures via: https://github.com/gillescastel/inkscape-figures
\usepackage{import}
\usepackage{pdfpages}
\usepackage{transparent}


% This command is the ONLY command imported via the packages snippet, think of it as a tiny package written out
\newcommand{\incfig}[2][1]{%
    \def\svgwidth{#1\columnwidth}
    \import{./figures/}{#2.pdf_tex}
}

% ╔═══════════════════════════════════════════════════════════╗
% ║                 bibliography and citing                   ║
% ╚═══════════════════════════════════════════════════════════╝

\usepackage{csquotes}
% \usepackage[backend=biber,citestyle=alphabetic,bibstyle=authortitle]{biblatex}
\usepackage{biblatex}

\addbibresource{bibliography.bib}

% ╔═══════════════════════════════════════════════════════════╗
% ║                    colours (colors)                       ║
% ╚═══════════════════════════════════════════════════════════╝

\definecolor{dkgreen}{rgb}{0,0.6,0}
\definecolor{gray}{rgb}{0.5,0.5,0.5}
\definecolor{mauve}{rgb}{0.58,0,0.82}
\definecolor{lightyellow}{rgb}{1,1,0.6}
\definecolor{reddish}{HTML}{A01A3A}

%  ╔══════════════════════════════════════════════════════════╗
%  ║                       book format                        ║
%  ╚══════════════════════════════════════════════════════════╝

\pagestyle{fancy}

\setlength\headheight{20pt}

\renewcommand{\sectionmark}[1]{\markright{\thesection.\ #1}}

\fancyfoot[C,CO]{\textcolor{gray}{Nathanael Chwojko-Srawley}}
\fancyfoot[R,RO]{\thepage}{}

% Create a new page style for the first page
\fancypagestyle{firstpage}{
  \fancyhf{} % Clear header and footer
  \fancyhead[L]{\theauthor}
  \fancyhead[R]{\today} % Date in the top-right
  \fancyfoot[C]{\textcolor{gray}{Nathanael Chwojko-Srawley}}
  \fancyfoot[R]{\thepage}
}

% Apply the first page style conditionally
\AtBeginDocument{\thispagestyle{firstpage}}

%have format the title command
\newcommand{\chapter}[1]{%
  \clearpage
  \vspace*{30pt} % Space before the chapter title
  {\normalfont\huge\bfseries #1\par} % Chapter title formatting
  \vspace{20pt} % Space after the chapter title
}


\setlength{\parindent}{0pt} % don't like indentation infront of paragraphs
\setlength{\parskip}{1ex}   %so that there is still space between paragarphs


%  ╔══════════════════════════════════════════════════════════╗
%  ║                       Shortcuts                          ║
%  ╚══════════════════════════════════════════════════════════╝

%\ExerciseLevelInToc

% I like original mathcal better, but need lowercase.
\DeclareFontFamily{OT1}{pzc}{}
\DeclareFontShape{OT1}{pzc}{m}{it}{<-> s * [1.10] pzcmi7t}{}
\DeclareMathAlphabet{\mathpzc}{OT1}{pzc}{m}{it}


% mathbb
\newcommand{\A}{{\mathbb{A}}}
\newcommand{\C}{{\mathbb{C}}}
\newcommand{\F}{{\mathbb{F}}}
\newcommand{\N}{{\mathbb{N}}}
\newcommand{\Q}{{\mathbb{Q}}}
\newcommand{\R}{{\mathbb{R}}}
\newcommand{\V}{{\mathbb{V}}}
\newcommand{\Z}{{\mathbb{Z}}}
\newcommand{\BI}{{\mathbb{I}}}
\renewcommand{\H}{{\mathbb{H}}}
\renewcommand{\P}{{\mathbb{P}}}


% mathcal
% \newcommand{\co}{{\mathfrak{o}}} %mathcal{o} looks like wreath product, so I changed it to mathfrak
\newcommand{\co}{{\mathpzc{o}}}
\newcommand{\CA}{{\mathcal{A}}}
\newcommand{\CB}{{\mathcal{B}}}
\newcommand{\CC}{{\mathcal{C}}}
\newcommand{\CD}{{\mathcal{D}}}
\newcommand{\CF}{{\mathcal{F}}}
\newcommand{\CG}{{\mathcal{G}}}
\newcommand{\CH}{{\mathcal{H}}}
\newcommand{\CI}{{\mathcal{I}}}
\newcommand{\CK}{{\mathcal{K}}}
\newcommand{\CL}{{\mathcal{L}}}
\newcommand{\CM}{{\mathcal{M}}}
\newcommand{\CN}{{\mathcal{N}}}
\newcommand{\CO}{{\mathcal{O}}}
\newcommand{\CQ}{{\mathcal{Q}}}
\newcommand{\CR}{{\mathcal{R}}}
\newcommand{\CS}{{\mathcal{S}}}
\newcommand{\CT}{{\mathcal{T}}}
\newcommand{\CV}{{\mathcal{V}}}
\newcommand{\CZ}{{\mathcal{Z}}}
% EXCPETION:
\newcommand{\CP}{{\mathcal{P}}}
% \newcommand{\CP}{\mathbb{C}\mathrm{P}}
\newcommand{\RP}{\mathbb{R}\mathrm{P}}


%Euscript
\newcommand{\EB}{\EuScript{B}}
\newcommand{\EC}{{\EuScript{C}}}
\newcommand{\EF}{{\EuScript{F}}}
\newcommand{\EL}{{\EuScript{L}}}
\newcommand{\EO}{{\EuScript{O}}}


%mathfrak (lowercase)
\newcommand{\fa}{{\mathfrak{a}}}
\newcommand{\fb}{{\mathfrak{b}}}
\newcommand{\fc}{{\mathfrak{c}}}
\newcommand{\fd}{{\mathfrak{d}}}
\newcommand{\fm}{{\mathfrak{m}}}
\newcommand{\fn}{{\mathfrak{n}}}
\newcommand{\fo}{{\mathfrak{o}}}
\newcommand{\fp}{{\mathfrak{p}}}
\newcommand{\fq}{{\mathfrak{q}}}


%mathfrak (upper case)
\newcommand{\FA}{{\mathfrak{A}}}
\newcommand{\FB}{{\mathfrak{B}}}
\newcommand{\FC}{{\mathfrak{C}}}
\newcommand{\FD}{{\mathfrak{D}}}
\newcommand{\FK}{{\mathfrak{K}}}
\newcommand{\FM}{{\mathfrak{M}}}
\newcommand{\FN}{{\mathfrak{N}}}
\newcommand{\FO}{{\mathfrak{O}}}
\newcommand{\FP}{{\mathfrak{P}}}


%mathscr
\newcommand{\ff}{{\mathscr{F}}}
\newcommand{\SA}{\mathscr{A}}
\newcommand{\SC}{\mathscr{C}}
\newcommand{\SF}{\mathscr{F}}
\newcommand{\SG}{\mathscr{G}}
\newcommand{\SH}{\mathscr{H}}
\newcommand{\SI}{\mathscr{I}}
\newcommand{\SK}{\mathscr{K}}
\newcommand{\SN}{\mathscr{N}}
\newcommand{\SQ}{\mathscr{Q}}
\newcommand{\SR}{\mathscr{R}}
\newcommand{\SV}{\mathscr{V}}
\newcommand{\SZ}{\mathscr{Z}}


% operatorname -- 2 letters (lowercase and upper case)
\newcommand{\ab}{{\operatorname{ab}}}
\newcommand{\ad}{{\operatorname{ad}}}
\newcommand{\ch}{{\operatorname{ch}}}
\newcommand{\ev}{{\operatorname{ev}}}
\newcommand{\id}{{\operatorname{id}}}
\newcommand{\im}{{\operatorname{im}}}
\newcommand{\nm}{{\operatorname{Nm}}}
\newcommand{\op}{{\operatorname{op}}}
\newcommand{\rk}{{\operatorname{rk}}}
\newcommand{\sh}{{\operatorname{sh}}}
\newcommand{\tr}{{\operatorname{tr}}}

\newcommand{\Ab}{{\operatorname{Ab}}}
\newcommand{\Br}{{\operatorname{Br}}}
\newcommand{\Ch}{{\operatorname{Ch}}}
\newcommand{\Cl}{{\operatorname{Cl}}}
\newcommand{\GL}{{\operatorname{GL}}}
\newcommand{\Nm}{{\operatorname{Nm}}}
\newcommand{\SL}{{\operatorname{SL}}}

\renewcommand{\Re}{{\operatorname{Re}}}
\renewcommand{\Im}{{\operatorname{Im}}}


%categories
\newcommand{\Grp}{{\textbf{Grp}}}
\newcommand{\Fld}{{\textbf{Fld}}}
\newcommand{\Sh}{{\textbf{Sh}}}
\newcommand{\Set}{{\textbf{Set}}}
\newcommand{\Mod}{{\textbf{Mod}}}
\newcommand{\PSh}{{\textbf{PSh}}}
\newcommand{\Sch}{{\textbf{Sch}}}
\newcommand{\Top}{{\textbf{Top}}}
\newcommand{\RgSp}{{\textbf{RgSp}}}
\newcommand{\Ring}{{\textbf{Ring}}}


%operatorname -- 3+ letters (lowercase and upper case)
\newcommand{\rad}{{\operatorname{rad}}}
\newcommand{\sep}{{\operatorname{sep}}}
\newcommand{\res}{{\operatorname{res}}}
\newcommand{\sgn}{{\operatorname{sgn}}}
\newcommand{\pre}{{\operatorname{pre}}}
\newcommand{\ord}{{\operatorname{ord}}}
\newcommand{\vol}{{\operatorname{vol}}}
\newcommand{\lcm}{{\operatorname{lcm}}}

\newcommand{\conj}{{\operatorname{conj}}}
\newcommand{\disc}{{\operatorname{disc}}}
\newcommand{\stab}{{\operatorname{stab}}}
\newcommand{\kdim}{{\operatorname{kdim}}}
\newcommand{\diag}{{\operatorname{diag}}}

\newcommand{\trdeg}{\operatorname{trdeg}}
\newcommand{\coker}{{\operatorname{coker}}}
\newcommand{\codim}{{\operatorname{codim}}}

\newcommand{\Ann}{{\operatorname{Ann}}}
\newcommand{\Aut}{{\operatorname{Aut}}}
\newcommand{\Der}{{\operatorname{Der}}}
\newcommand{\Div}{{\operatorname{Div}}}
\newcommand{\Emb}{{\operatorname{Emb}}}
\newcommand{\End}{{\operatorname{End}}}
\newcommand{\Ext}{{\operatorname{Ext}}}
\newcommand{\Gal}{{\operatorname{Gal}}}
\newcommand{\Hom}{{\operatorname{Hom}}}
\newcommand{\Ind}{{\operatorname{Ind}}}
\newcommand{\Inn}{{\operatorname{Inn}}}
\newcommand{\Int}{{\operatorname{int}}}
\newcommand{\Irr}{{\operatorname{Irr}}}
\newcommand{\Jac}{{\operatorname{Jac}}}
\newcommand{\Mor}[1]{\operatorname{Mor}(\textbf{#1})}
\newcommand{\Nil}{{\operatorname{Nil}}}
\newcommand{\Obj}{{\operatorname{Obj}}}
\newcommand{\Out}{{\operatorname{Out}}}
\newcommand{\Pic}{{\operatorname{Pic}\,}}
\newcommand{\Rep}{{\operatorname{Rep}}}
\newcommand{\Res}{{\operatorname{Res}}}
\newcommand{\Spl}{{\operatorname{Spl}}}
\newcommand{\Syl}{{\operatorname{Syl}}}
\newcommand{\Sym}{{\operatorname{Sym}}}
\newcommand{\Tor}{{\operatorname{Tor}}}

\newcommand{\Char}{{\operatorname{char}}}
\newcommand{\Frac}{{\operatorname{Frac}}}
\newcommand{\Frob}{{\operatorname{Frob}}}
\newcommand{\Proj}{{\operatorname{Proj}\,}}
\newcommand{\RMod}{{}_R\operatorname{Mod}}
\newcommand{\SExt}{{\mathcal{E}\emph{xt}}}
\newcommand{\SHom}{{\emph{Hom}}}
\newcommand{\Span}{{\operatorname{span}}}
\newcommand{\Spec}{{\operatorname{Spec}\,}}
\newcommand{\Supp}{{\operatorname{Supp}}}
\newcommand{\Tens}[1]{#1\otimes --}
\newcommand{\fHom}[1]{\operatorname{Hom}(#1, --)}

\newcommand{\mSpec}{{\operatorname{mSpec}}}
\newcommand{\cofHom}[1]{\operatorname{Hom}(--, #1)}


% connectors
% \renewcommand{\mid}{\big|}
\newcommand{\sse}{\subseteq}
\newcommand{\subg}{\leqslant}
\newcommand{\supg}{\geqslant}
\newcommand{\ssne}{\subsetneq}
\newcommand{\isosubg}{\lesssim}
\newcommand{\nsse}{\not\subseteq}
\newcommand{\nsubg}{\trianglelefteq}
\newcommand{\nsupg}{\trianglerighteq}
\newcommand{\csubg}{\blacktriangleleft}
\renewcommand{\mid}{\big|}


% Arrows
\newcommand{\rw}{\rightarrow}
\newcommand{\Rw}{\Rightarrow}
\newcommand{\lw}{\leftarrow}
\newcommand{\Lw}{\Leftarrow}
\newcommand{\lrw}{\leftrightarrow}
\newcommand{\LRw}{\Leftrightarrow}
\newcommand{\acton}{\curvearrowright}
\newcommand{\actson}{\curvearrowright}
\newcommand\mapsfrom{\mathrel{\reflectbox{\ensuremath{\mapsto}}}}

% arrows for commutative diagram: create four new angled double-struck arrows
\newcommand\myrot[1]{\mathrel{\rotatebox[origin=c]{#1}{$\Rightarrow$}}}
\newcommand\NEarrow{\myrot{45}}
\newcommand\NWarrow{\myrot{135}}
\newcommand\SWarrow{\myrot{-135}}
\newcommand\SEarrow{\myrot{-45}}


%shorthands
\newcommand{\oln}[1]{{\overline{#1}}}
\renewcommand{\bf}[1]{\textbf{#1}}
\newcommand{\qn}{\quad\newline}


% limits
\newcommand{\Lim}{{\varprojlim}}
\newcommand{\coLim}{{\varinjlim}}
\newcommand{\Dlim}{\lim\limits_{\longrightarrow}}
\newcommand{\coDlim}{\lim\limits_{\longleftarrow}}


%for saturated ideals in the localization section (I don't like these, I think I'll delete)
\newcommand{\LIm}{{}^e }
\newcommand{\LPre}{{}^c }

%  ╔══════════════════════════════════════════════════════════╗
%  ║                    Custom Commands                       ║
%  ╚══════════════════════════════════════════════════════════╝

% swap varphi and phi
\let\temp\phi
\let\phi\varphi
\let\varphi\temp

\newcommand{\placeholder}{\_\_\_\_\_}

\newcommand{\set}[2]{\left\{#1 \ \middle|\ #2\right\}}

%mod should not have the crazy amount of space to its right!
\renewcommand{\mod}{\,\operatorname{mod}\,}

% dcups
\newcommand{\dcup}{\sqcup}
\newcommand{\Dcup}{\coprod}
% \newcommand{\Dcup}{\bigsqcup}

\newcommand{\nbot}{\mathrel{\rlap{$\bot$}\mkern2mu/}}

%a nice large circle
\newcommand{\tikzcircle}[2][red,fill=black]{\tikz[baseline=-0.5ex]\draw[#1,radius=#2] (0,0) circle ;}%

\makeatletter
\newcommand*\bigcdot{\mathpalette\bigcdot@{.5}}
\newcommand*\bigcdot@[2]{\mathbin{\vcenter{\hbox{\scalebox{#2}{$\m@th#1\bullet$}}}}}
\makeatother

%somehow not a default command
\def\rddots{\mathstrut^{.^{.^{.}}}}

%% new delimiter
\DeclarePairedDelimiter{\ceil}{\lceil}{\rceil}


%  ╔══════════════════════════════════════════════════════════╗
%  ║                     environments                         ║
%  ╚══════════════════════════════════════════════════════════╝


%remember the option: breakable

% create tcolor boxes
\newtcbtheorem[number within=section]{thm}{Theorem}%
{
colback=green!5,
colframe=green!35!black,
fonttitle=\bfseries,
label type=theorem
}{th}

\newtcbtheorem[number within=section, use counter from=thm]{lem}{Lemma}%
{
colback=blue!5,
colframe=black!35!black,
fonttitle=\bfseries,
label type=lemma
}{lm}
\newtcbtheorem[number within=section, use counter from=thm]{prop}{Proposition}%
{
colback=white!5,
colframe=white!35!black,
fonttitle=\bfseries,
label type=proposition
}{pr}
\newtcbtheorem[number within=section, use counter from=thm]{cor}{Corollary}%
{colback=blue!5,
colframe=blue!35!black,
fonttitle=\bfseries,
label type=corollary
}{co}
\newtcbtheorem[number within=section, use counter from=thm]{axiom}{Axiom}%
{colback=black!5,
colframe=yellow!35!black,
fonttitle=\bfseries,
label type=axiom
}{ax}
\newtcbtheorem[number within=section, use counter from=thm]{defn}{Definition}%
{colback=red!5,
colframe=red!35!black,
fonttitle=\bfseries,
label type=definition
}{df}


% for <s-cr> to create \item[$\LRw$] instead of \item
\newenvironment{equivEnumerate}
 {\begin{enumerate}
	}{
\end{enumerate}}



%Custom enviroments
 \newenvironment{note}
 {\begin{tcolorbox}[enhanced, sharp corners, colback=white , borderline={0.2pt}{0pt}{black}]
   \begin{quote}\textbf{Note}:
   }{
   \end{quote}\end{tcolorbox}}

 \newenvironment{intuition}
 {\begin{quote}\textbf{Intuition}:
   }{
   \end{quote}}

\newtcbtheorem[number within=section, use counter from=thm]{titledBox}{}{%
  enhanced,
  sharp corners,
  colback=white,
  colframe=black,
  borderline={0.2pt}{0pt}{black},
  left=2mm,
  right=2mm,
  top=2mm,
  bottom=2mm,
  title style={opacity=1},
  attach title to upper,
  before upper={\textbf{\tcbtitletext}\quad},
  coltitle=black,
  fonttitle=\bfseries,
  separator sign={\quad}
}{box}


%better example and proof environments
\def\exampletext{Example} % If English
\colorlet{colexam}{red!55!black} % Global example color


\newtcbtheorem[number within=section, use counter from=thm]{example}{Example}{% \newtcolorbox[use counter=testexample]{testexamplebox}{%
% Example Frame Start
empty,% Empty previously set parameters
title={\exampletext \thetcbcounter #1},% use \thetcbcounter to access the testexample counter text
% Attaching a box requires an overlay
attach boxed title to top left,
   % Ensures proper line breaking in longer titles
   minipage boxed title,
% (boxed title style requires an overlay)
boxed title style={empty,size=minimal,toprule=0pt,top=4pt,left=3mm,overlay={}},
coltitle=colexam,fonttitle=\bfseries,
before=\par\medskip\noindent,parbox=false,boxsep=0pt,left=3mm,right=0mm,top=2pt,breakable,pad at break=0mm,
   before upper=\csname @totalleftmargin\endcsname0pt, % Use instead of parbox=true. This ensures parskip is inherited by box.
% Handles box when it exists on one page only
overlay unbroken={\draw[colexam,line width=.5pt] ([xshift=-0pt]title.north west) -- ([xshift=-0pt]frame.south west); },
% Handles multipage box: first page
overlay first={\draw[colexam,line width=.5pt] ([xshift=-0pt]title.north west) -- ([xshift=-0pt]frame.south west); },
% Handles multipage box: middle page
overlay middle={\draw[colexam,line width=.5pt] ([xshift=-0pt]frame.north west) -- ([xshift=-0pt]frame.south west); },
% Handles multipage box: last page
overlay last={\draw[colexam,line width=.5pt] ([xshift=-0pt]frame.north west) -- ([xshift=-0pt]frame.south west); }%
}{ex}



\def\prooftext{Proof} % If English
\NewDocumentEnvironment{Proof}{ O{} }
{
    \colorlet{colexam}{black!55!black} % Global example color
    \newtcolorbox[number within=section]{testproofbox}{%
        empty,% Empty previously set parameters
        title={\emph{\prooftext} \thetcbcounter: #1},
        attach boxed title to top left,
        minipage boxed title,
        boxed title style={empty,size=minimal,toprule=0pt,top=4pt,left=3mm,overlay={}},
        coltitle=colexam,
        fonttitle=\bfseries,
        before=\par\medskip\noindent,
        parbox=false,
        boxsep=0pt,
        left=3mm,
        right=0mm,
        top=2pt,
        breakable,
        pad at break=0mm,
        before upper=\csname @totalleftmargin\endcsname0pt,
        % Removed height constraints
        % Added flexible width
        width=\dimexpr\textwidth-3mm\relax,
        % Modified overlays
        overlay unbroken={\draw[colexam,line width=.5pt] ([xshift=-0pt]title.north west) -- ([xshift=-0pt]frame.south west); },
        overlay first={\draw[colexam,line width=.5pt] ([xshift=-0pt]title.north west) -- ([xshift=-0pt]frame.south west); },
        overlay middle={\draw[colexam,line width=.5pt] ([xshift=-0pt]frame.north west) -- ([xshift=-0pt]frame.south west); },
        overlay last={\draw[colexam,line width=.5pt] ([xshift=-0pt]frame.north west) -- ([xshift=-0pt]frame.south west); },%
    }
    \begin{testproofbox}
}
{\end{testproofbox}\endlist}


%  ╔══════════════════════════════════════════════════════════╗
%  ║                    for Dynkin diagrams                   ║
%  ╚══════════════════════════════════════════════════════════╝

\def\row#1/#2!{#1_{\IfStrEq{#2}{}{n}{#2}} & \dynkin{#1}{#2}\\}
\newcommand{\tble}[1]{
   \renewcommand*\do[1]{\row##1!}
   \[
      \begin{array}{ll}\docsvlist{#1}\end{array}
   \]
}

%  ╔══════════════════════════════════════════════════════════╗
%  ║                         links                            ║
%  ╚══════════════════════════════════════════════════════════╝

\definecolor{LinkColour}{HTML}{A64A3B} % favourite
% \definecolor{LinkColour}{HTML}{8C3D32} % 2nd place
\hypersetup{
  colorlinks,
  citecolor=black,
  filecolor=magenta,
  linkcolor=LinkColour,
  urlcolor=blue,
}


%  ╔══════════════════════════════════════════════════════════╗
%  ║                          misc                            ║
%  ╚══════════════════════════════════════════════════════════╝

\stackMath %to be able to stack text in math mode
\makeindex

%import options: all, bibliography, book-formatting, booksSetting, customEnvironments, dynkin, hw-formatting, language-setup, newcommands, packages, preamble, proofs, settings, tcolorbox, test.svg,


\begin{document}
\maketitle

The blog was inspired by the following thought process:
\begin{quoting}
    We start with the natural numbers, construct the integers, then the rationals. We complete them to get
    the reals, and finally adjoin $i$ to get the complex numbers. This is the ``standard'' way. But is it the \emph{only} way?
\end{quoting}

Imagine an alien civilization that understands number theory perfectly but has no physical intuition for
``continuity'' or ``measuring lengths.'' To them, the algebraic closure of the rationals, $\oln{\Q}$, is the
natural universe. This blog explores how such a civilization would discover the real numbers not as
a foundation, but as a strange, derived sub-structure, and how this would change their understanding of
calculus.

Be warned that this blog is written with ``free-flow'' writing and not too much thought was put in to the
structure. I will not really prove anything and not define all my terms; I will just assume proficiency with
most of the background material. All the theorems and relevant definitions can be found in
either~\cite{NateAlgebra}, \cite{NateAlgGeo}, \cite{NateTop}, \cite{NateComplexAnalysis}, or \cite{NateCFT}.

\section{Constructing Reality from Algebra}

Our aliens start with $\Q$. Their natural next step is to solve all polynomial equations, leading them to
$\oln{\Q}$. They have no topology yet. How do they find $\R$? Classically, we view $\C$ as a degree 2
extension of $\R$. The aliens view this in reverse: they search for a subfield of index 2 inside their
algebraic closure.

\begin{thm}{Artin-Schreier}{artin}
Let $K$ be an algebraically closed field. If $F \subset K$ is a subfield of finite index (i.e., $1 < [K:F] < \infty$), then:
\begin{enumerate}
    \item $[K:F] = 2$.
    \item $K = F(i)$ where $i^2 = -1$.
    \item $F$ is a \emph{Real Closed Field}.
\end{enumerate}
\end{thm}

To find ``Reality,'' the aliens look at the symmetries of their universe: the Absolute Galois Group $G = \Gal(\oln{\Q}/\Q)$. They seek an \emph{involution} $\sigma \in G$ (an automorphism where $\sigma^2 = \id$).

\begin{defn}{The Alien Construction of $\R$}{alienreal}
    Let $\sigma \in \Gal(\oln{\Q}/\Q)$ be an involution. We define the ``Real Algebraic Numbers'' relative to $\sigma$ as the fixed field:
    \[ F_\sigma = \{ x \in \oln{\Q} \mid \sigma(x) = x \} \]
\end{defn}

\begin{note}
    This construction is \emph{not unique}. There are infinitely many involutions in the Galois group, all conjugate to each other. The alien can build \emph{a} real line, but they cannot distinguish \emph{which} one is ``ours'' without making an arbitrary choice. To them, $\sqrt[3]{2}$ might be ``real'' in one coordinate system, but ``complex'' in another.
\end{note}

\subsection{The Topological Completion}
Once $F_\sigma$ is chosen, the ordering is algebraic: $x \ge 0 \iff \exists y \in F_\sigma, x=y^2$. But how do they recover $\C$? They cannot use a topology intrinsic to $\oln{\Q}$ alone. Instead, they define a \emph{norm} relative to their chosen involution.

Since $F_\sigma$ is real closed, the sum of two squares is always a square. The alien defines the ``$\sigma$-norm'' $|\cdot|_\sigma : \oln{\Q} \to F_\sigma$ as:
\[ |z|_\sigma = \sqrt{z \cdot \sigma(z)} \]
Algebraically, if $z = x + iy$, then $z \cdot \sigma(z) = x^2 + y^2$, which has a square root in $F_\sigma$. This allows them to define a metric $d(a, b) = |a - b|_\sigma$.

They now run the standard completion process (Cauchy sequences):
\begin{enumerate}
    \item Completing $F_\sigma$ yields $\R$.
    \item Completing $\oln{\Q}$ (which is a 2D vector space over $F_\sigma$) yields the product topology:
    \[ \widehat{\oln{\Q}} \cong \widehat{F_\sigma} \oplus \widehat{F_\sigma} i \cong \R \oplus \R i \cong \C \]
\end{enumerate}

\section{The Archimedean Anomaly}

The topology derived above---the Order Topology---would be viewed by the aliens as algebraically natural, but geometrically ``dangerous.''

\subsection{The Topology of Squares}
The aliens admit the definition is elegant because it requires no new concepts. They don't need a visualization of a number line; they just use squares.
Rather than defining open intervals $(a,b)$ via inequalities, they define the basis of the topology as:
\[
U_{a,b} = \{ x \in F \mid (x-a) \in (F^\times)^2 \text{ and } (b-x) \in (F^\times)^2 \}
\]
Because $F$ is a Real Closed Field, the ``positive'' elements are exactly the squares. So, the alien sees the
order topology not as ``left vs. right,'' but as a statement about the solvability of quadratic equations. To be ``in the interval'' means you can take the square root of your distance to the endpoints.

\subsection{The Singularity}
However, once they start studying this topology, they find it completely bizarre compared to the rest of their universe ($\Q_p$, finite fields).

Most valuations in algebra are \textbf{Non-Archimedean} (satisfying the strong triangle inequality $|x+y| \le \max(|x|, |y|)$). This creates a stable, bounded universe. The Order Topology on $F$ is the only one that is \textbf{Archimedean} ($|x+y| \le |x| + |y|$).

\begin{itemize}
    \item \textbf{Explosive Growth}: $F$ is the only place where adding small things together eventually creates an arbitrarily large monster. In $\Q_p$, ``infinity'' is not a direction you can reach by walking. In $F$, it is.
    \item \textbf{Directedness}: The existence of the order creates an asymmetry. In $\C$ or $\Q_p$, $x$ and $-x$ are topologically indistinguishable. In $F$, there is a fundamental difference between ``positive'' (a square) and ``negative'' (not a square).
\end{itemize}

To the alien, the universe is mostly disconnected dust ($\Q_p$), but there is one \textbf{singular point at infinity} ($\R$) where the dust fuses together, explodes in size, and gains a direction.

\section{The Rigidity Paradox}

The standard modern intuition is to consider real calculus (smooth functions, bump functions) to be the standard, and complex calculus (holomorphic functions) to be a restrictive special case satisfying the Cauchy-Riemann equations.

To the alien, this is backward. The ``Natural'' calculus is on $\oln{\Q}$ (or $\C$): a function $f: \C \to \C$
is differentiable if it is locally linear \emph{over the field} $\C$. This means the local approximation must
be multiplication by a scalar $\lambda \in \C$.

If we translate this requirement into the language of the fixed field $\R \subset \C$, the requirement that the Jacobian matrix acts like a complex number \emph{is} the Cauchy-Riemann equations.

\[
\begin{pmatrix} a & -b \\ b & a \end{pmatrix} \iff \text{Cauchy-Riemann}
\]

There is another reason the alien considers Cauchy-Riemann functions the only ``natural'' ones. Remember, they constructed their coordinates by making an \emph{arbitrary choice} of $\sigma$.

A ``natural'' function (like a polynomial $P(z) = z^2+1$) is defined using only field operations; it does not care which $\sigma$ you chose. A ``non-natural'' function (like $f(z) = \Re(z) = \frac{z+\sigma(z)}{2}$) requires you to know exactly which involution was chosen.

The Cauchy-Riemann equations can be written in Wirtinger derivative notation as $\frac{\partial f}{\partial \bar{z}} = 0$. Translated to Alien-speak:
\[ \frac{\partial f}{\partial z^\sigma} = 0 \]
To the alien, the Cauchy-Riemann equations are the statement: \emph{``This function depends only on the number $z$ itself, and not on the arbitrary choice of involution used to construct the coordinate system.''}

\subsection{Why is Real Calculus ``Freer''?}
If $\C$ is the natural world, why does $\R$ allow for ``wild'' objects like bump functions (smooth functions that are non-zero in a small set and zero everywhere else)?

\begin{itemize}
    \item In $\C$: To be smooth, you must satisfy the Cauchy-Riemann equations (Coordinate Independence). This ``cross-check'' enforces rigidity: if you know a function on a tiny open set, you know it everywhere (Analytic Continuation).
    \item In $\R$: By restricting to the fixed field, we effectively \emph{turn off} the sensor in the imaginary dimension. We no longer require the derivative to match an approach from the $i$-direction.
\end{itemize}

The aliens view Smooth Real functions not as a richer class, but as a \emph{broken} class. They are the pathological consequence of taking a slice of the universe and ignoring the ambient constraints of the algebraic closure.

\section{Algebraic Differentiation: The Dual Numbers}

Can the aliens define differentiation without limits or topology? Yes. They use the \emph{Dual Numbers}.

\begin{defn}{Algebraic Derivative}{algder}
    Let $R$ be a ring. Consider the ring of dual numbers $R[\varepsilon]/(\varepsilon^2)$. For a polynomial $P(x)$, we calculate:
    \[ P(x + \varepsilon) = P(x) + P'(x)\varepsilon \]
    The coefficient of $\varepsilon$ is the derivative.
\end{defn}

This definition is purely algebraic. Does it descend to the Real subfield?

\begin{center}
\begin{tikzcd}
    K[x] \arrow[r, "D"] \arrow[d, "\text{restrict}"'] & K[x] \arrow[d, "\text{restrict}"] \\
    F[x] \arrow[r, "D"] & F[x]
\end{tikzcd}
\end{center}

Since the derivative of a polynomial with coefficients in $F$ has coefficients in $F$, the diagram commutes.

What is interesting about using automatic differentiation is that the input is forced to be a algebraic
(polynomials, rational polynomials, series, etc.). The Algebraic Derivative is blind to the ``smooth
monsters'' of real analysis. If you try to feed it a bump function (which has no Taylor series at the
boundary), the algebraic machine jams. To the alien, bump functions are ghosts that appear only when you
introduce topology.

\section{The Spectrum of Topologies}

Finally, let's take a look at Valuation Theory. They know that for every prime number $p$ (and the infinite
prime $\infty$), there is a valid, natural metric topology on $\oln{\Q}$. This gives them a whole spectrum of
possible universes.

Which one is the ``correct'' one? They run a simple compatibility test against their constructed Real Field $F_\sigma$.

\begin{itemize}
    \item \textbf{The $p$-adic Topology}: If they put a $p$-adic metric on $\oln{\Q}$, the subfield $F_\sigma$ becomes totally disconnected dust. The algebraic order ($x > y$) has no relation to the topological closeness. The interval $[0,1]$ is shredded.
    \item \textbf{The Archimedean Topology}: If they put the infinite-prime metric on $\oln{\Q}$, the subfield $F_\sigma$ becomes a connected line. The topology respects the order.
\end{itemize}

The alien concludes that while $\oln{\Q}$ has many natural topologies, only the Archimedean one is compatible with the Order Structure of the fixed field.

\begin{example}{The Local Constancy Failure}{locallyConstantFails}
    Differentiation is a universal operator, but its meaning changes with the topology. In the $p$-adic worlds ($\Q_p$), the topology is totally disconnected. There exist locally constant functions that are not globally constant.
    \[ f'(x) = 0 \not\Rw f(x) = \text{const} \]
    In $\Q_\infty$ ($\R$), the Archimedean property forces the topology to be connected.
    \[ f'(x) = 0 \implies f(x) = \text{const} \]
\end{example}

The aliens conclude that $\R$ is not special just because of its order, but because it is the \emph{Connected Field}. It is the only completion of $\Q$ where the ``micro-behavior'' (derivative) forces the ``macro-behavior'' (continuity) to align.

This connectivity ascends to the algebraic closure, $\C$. Because the topology of $\C$ is just the product
topology of $\R$, $\C$ inherits this cohesion. This is the structural reason why \emph{Analytic Continuation}
works in the Archimedean world (a function defined on a small disk determines the function everywhere) but
fails in the $p$-adic worlds. The ``rigid'''' structure of $\C$ is induced by the real order
Topology.

\printbibliography

\end{document}
